% Mathematical section. Citation keys repoObservable and repoMeasureTransport
% should identify the currently fetched November 20 ProveIt packets.
\section{A common probability space and its represented completions}
\label{sec:prob-foundation}

The coherent-partition and completion interfaces in the current ProveIt
framework already provide the appropriate starting point
\cite{repoObservable,repoMeasureTransport}. We specialize them to probability,
add an explicit distinction between arbitrary random variables and integrable
ones, and spell out the certificates needed by discontinuous event operations.
These are constructions on one common probability space. A list of laws with
unrelated pairwise couplings would not supply the joint structure required by
addition, independence, conditioning, or events involving several variables.

\subsection{The exact finite base}

For $n\in\mathbb N$, let $\Omega_n=\{0,1\}^n$, coded by the integers below
$2^n$. Each word has mass $2^{-n}$. A cylinder event is a Boolean table on
$\Omega_n$; a rational cylinder variable is a rational table on $\Omega_n$.
When $m\ge n$, lift a table to $\Omega_m$ by ignoring the last $m-n$ bits.
Every row is then repeated $2^{m-n}$ times. Two or more tables are compared
after lifting them to the largest of their levels. The projection from level
$l$ to $n$ is the composite of the projections through any intermediate
level $m$, so these joint refinements are coherent.

For finite tables write
\[
 \mathbb P_n(A)=2^{-n}\sum_{\omega\in\Omega_n}\mathbf1_A(\omega),
 \qquad I_n(f)=2^{-n}\sum_{\omega\in\Omega_n}f(\omega).
\]
Repetition preserves both quantities. Thus the subscript may be omitted once
a common level is understood. Every construction here is primitive recursive:
the number of rows may be very large, but there is an explicit finite bound.
Rationals are signed integers over positive denominators; comparison is exact
cross multiplication. The denominator guard is retained in every formula.

\begin{lemma}[The finite probability kernel]\label{lem:prob-finite}
The cylinder algebra has exact finite additivity, monotonicity, complement
and union laws. Rational integration is positive and linear, and
\[
 |I(f)-I(g)|\le I(|f-g|),\qquad
 I(|f+g|)\le I(|f|)+I(|g|).
\]
If $\varphi$ is rationally evaluated and $L$-Lipschitz, then
$I(|\varphi(f)-\varphi(g)|)\le L I(|f-g|)$.
For finite tables on independent blocks of bits,
\[
 I(f(\omega)g(\eta))=I(f)I(g).
\]
\end{lemma}
\begin{proof}
The event statements follow from Boolean identities on each row. For example,
$\mathbf1_{A\cup B}+\mathbf1_{A\cap B}=\mathbf1_A+\mathbf1_B$.
The integral inequalities follow by multiplying the corresponding rational
pointwise inequalities by the nonnegative row mass and summing. For the last
identity, write the expectation as a finite double sum and use distributivity:
$\sum_{i,j}f_i g_j=(\sum_i f_i)(\sum_j g_j)$. Normalizing by the product of
the two word counts gives the assertion.

For the iPRA proof, clear the positive denominators and use the partial-sum
inequality as the induction invariant. It is quantifier-free. Row conditions
are decidable Boolean tests; their bounded conjunction is encoded by a Boolean
fold. Universal presentation hypotheses are instantiated at the finitely many
rows used. No induction on a statement about all precisions is needed.
\end{proof}

Natural-weight finite probability spaces are also available: replace
$2^{-n}$ by $w_i/W$, where $W=\sum_iw_i>0$. They are useful for rational
quadrature laws. Their mass vectors need not be exactly dyadic, so they should
not be silently identified with a single finite fair-bit table. They can be
approximated by dyadic masses with explicit rounding error, or retained as
finite auxiliary laws with separately supplied transports.

\subsection{Three different distances}

On rational cylinder variables define
\[
 d_0(f,g)=I(\min(1,|f-g|)),\qquad
 d_1(f,g)=I(|f-g|),\qquad
 d_2(f,g)^2=I((f-g)^2).
\]
The first two are exactly rational. Square roots in $d_2$ may instead be
handled by squared rational inequalities, or by rational bisection with a
proved error bound. On cylinder events define
\[
 d_{\rm ev}(A,B)=\mathbb P(A\mathbin{\triangle}B)
                =d_1(\mathbf1_A,\mathbf1_B).
\]
The triangle inequality for $d_0$ follows from
\[
 \min(1,|x-z|)\le\min(1,|x-y|)+\min(1,|y-z|).
\]
The triangle inequality for $d_1$ is Lemma~\ref{lem:prob-finite}. The usual
$d_2$ inequality follows from finite Cauchy--Schwarz; its square-free kernel is
\[
 \left(\sum_i w_i a_i b_i\right)^2
 \le\left(\sum_iw_i a_i^2\right)\left(\sum_iw_i b_i^2\right),
\]
proved by expanding the nonnegative sum
$\sum_{i,j}w_iw_j(a_i b_j-a_j b_i)^2$.

\begin{lemma}[Finite probability interpretation of $d_0$]
For rational $0<\delta\le1$,
\[
 \delta\,\mathbb P(|f-g|>\delta)
 \le d_0(f,g)
 \le \delta+\mathbb P(|f-g|>\delta).
\]
\end{lemma}
\begin{proof}
On a row with $|f-g|>\delta$, the capped difference is at least $\delta$;
on every other row it is at most $\delta$. On exceptional rows it is at most
one. Sum these two elementary inequalities.
\end{proof}

Thus $d_0$ controls convergence in probability without requiring the first
absolute moment. This distinction is essential for continuous laws such as
the Cauchy distribution. An $L^1$-only foundation would exclude its identity
random variable even though its distribution and all bounded observables are
well defined.

\subsection{Names, equality, and their arithmetic meaning}

Put $\varepsilon_k=2^{-k}$. For any one of the displayed metrics, a
\emph{regular name} is a fixed primitive-recursive family $x_k$ of finite
tables, with their levels, satisfying
\begin{equation}\label{eq:prob-regular}
 d(x_k,x_l)\le\varepsilon_k+\varepsilon_l\qquad(k,l\in\mathbb N).
\end{equation}
Its limit notation denotes a represented object. Two names are equal when
\begin{equation}\label{eq:prob-equal}
 d(x_k,y_l)\le\varepsilon_k+\varepsilon_l\qquad(k,l\in\mathbb N).
\end{equation}
This is an arithmetical all-precision relation, not a Boolean equality test.
It needs no quotient type inside iPRA.

For completeness, the small rational fact used by this definition is worth
recording. If rational $a,b$ satisfy $a\le b+2^{-m}$ for all $m$, then $a\le b$.
Decide their rational comparison. If $a>b$ and $D$ is a denominator of their
difference, then $a-b\ge1/D$; choosing $m=D+1$ contradicts the displayed
inequality because $2^{D+1}>D$. This is a primitive-recursively bounded
rational argument. Applying it to a triangle inequality with the intermediate
name evaluated at precision $m$ proves transitivity of
\eqref{eq:prob-equal}. Reflexivity and symmetry are immediate.

An $L^0$ random variable means a $d_0$ name. An integrable random variable is
a $d_1$ name; a square-integrable random variable is a $d_2$ name. Their
forgetful maps are justified by $d_0\le d_1\le d_2$, the second inequality
being finite Cauchy--Schwarz with total mass one. An event means a
$d_{\rm ev}$ name of Boolean tables. It is identified modulo null distance;
it does not carry a general procedure deciding whether an individual infinite
bit sequence belongs to it.

These names form the represented part of the usual external metric
completions. They are closed under supplied jointly PR Cauchy families with
explicit PR moduli, by a diagonal precision choice. They are not claimed to
form a classically complete space under arbitrary Cauchy sequences. The
ordinary measure-theoretic interpretation is an adequacy statement; it is
not an additional axiom of the internal arithmetic.

\begin{proposition}[Continuous finite operations extend]
Suppose a fixed PR approximate table operation $F_p$ satisfies
\[
 d(F_p(a),F_q(b))\le Ld(a,b)+\varepsilon_p+\varepsilon_q
\]
for a supplied rational $L\ge0$. If $L\le2^h$, then
$y_k=F_{k+2}(x_{k+h+2})$ is a regular output name, depending only on the
represented input.
\end{proposition}
\begin{proof}
Substituting \eqref{eq:prob-regular} makes the right-hand side at most
$(\varepsilon_k+\varepsilon_l)/2$. Substituting
\eqref{eq:prob-equal} gives the same bound for equivalent inputs. All
precision choices are explicit PR expressions. This is the common completion
rule already used in the repository, now applied to probability tables.
\end{proof}

Consequently addition, rational scaling, absolute value, positive and negative
parts, minimum, maximum, and clipping have their expected representations.
For example addition uses one guard bit because the two errors add. A
bounded Lipschitz test function gives a real expectation from every $L^0$
name: if $|h|\le M$ and $h$ has Lipschitz constant $L$, then
\[
 I(|h(f)-h(g)|)\le \max(L,2M)d_0(f,g).
\]
Indeed, for $|f-g|\le1$ use Lipschitz continuity; otherwise use the bound
$2M$. This provides the law of a general random variable through bounded
observables even when the identity has no expectation.

Every $L^0$ name also supplies effective tightness. Choose one finite
approximation $x_r$ and a rational $B$ bounding all its values. For every
later approximation $x_s$,
\[
 \mathbb P(|x_s|>B+1)\le d_0(x_s,x_r)
                         \le\varepsilon_s+\varepsilon_r.
\]
Thus choosing $r$ from the desired error gives an explicit tail radius.
Effective continuous maps with supplied PR uniform moduli on rational compact
intervals extend to $L^0$: first discard this small tail, then use the local
modulus on the compact part. A reciprocal or logarithm additionally needs
an explicit bound on the probability of approaching its singular set.

\subsection{Integration, events, and the threshold issue}

\begin{theorem}[Represented expectation]
If $X$ has an $L^1$ name $x_k$, then $I(x_k)$ is a regular rational real
name. It defines $\mathbb E X$, independent of the name, and
\[
 |\mathbb E X-\mathbb E Y|\le\|X-Y\|_1.
\]
The resulting expectation is linear and positive. The completed order is
$X\le Y$ when $\|(X-Y)^+\|_1=0$.
\end{theorem}
\begin{proof}
Apply the exact finite inequality
$|I(x_k)-I(x_l)|\le d_1(x_k,x_l)$ and then
\eqref{eq:prob-regular}. Equivalent names give equivalent scalar names by
the same inequality. Every linear identity holds exactly on a common finite
table and passes to names with its explicitly budgeted errors. Positivity
follows because the negative part has zero norm. This uses order on represented
classes, not a decision about pointwise order of completed values.
\end{proof}

\begin{proposition}[Represented Boolean event algebra]
Complement is an isometry for $d_{\rm ev}$, and union and intersection satisfy
\[
 d_{\rm ev}(A\cup B,A'\cup B'),\quad
 d_{\rm ev}(A\cap B,A'\cap B')
 \le d_{\rm ev}(A,A')+d_{\rm ev}(B,B').
\]
They extend to represented events, preserving all Boolean identities modulo
null distance. Probability is $\mathbb P(A)=\mathbb E\mathbf1_A$ and satisfies
the finite probability laws at every precision.
\end{proposition}
\begin{proof}
Changing either Boolean operation can only affect a row where at least one
input indicator changes. The finite union bound proves the estimate.
Complement uses the same level; the binary operations use one guard bit.
Their exact finite identities and the expectation contraction then prove the
asserted completed identities.
\end{proof}

There is an essential qualification. The map $x\mapsto\mathbf1_{x\le t}$
is discontinuous. Small $L^0$ or $L^1$ distance alone does not make its
threshold indicators close. For instance the constant tables $1/n$ converge
to zero, while $\mathbf1_{1/n\le0}$ is always zero and
$\mathbf1_{0\le0}$ is one.

\begin{lemma}[Threshold lifting with a boundary certificate]
Let $x_r$ be an $L^0$ name and let $t$ be rational. Suppose PR functions
$b(k),R(k)$ satisfy, at every $r\ge R(k)$,
\[
 \mathbb P(|x_r-t|\le2^{-b(k)})\le2^{-k}.
\]
Then the threshold events $\{x_r\le t\}$ can be reindexed by a PR function
to form a regular represented event.
\end{lemma}
\begin{proof}
For two finite tables $x_r,x_s$ and $0<\delta\le1$, a change of threshold
indicator implies either $|x_r-t|\le\delta$ or $|x_r-x_s|>\delta$.
Consequently
\[
 d_{\rm ev}(\{x_r\le t\},\{x_s\le t\})
 \le\mathbb P(|x_r-t|\le\delta)+\delta^{-1}d_0(x_r,x_s).
\]
At requested precision $k$, choose $\delta=2^{-b(k+3)}$ and choose a
monotone precision schedule $r(k)$ at least
$R(k+3)$ and $b(k+3)+k+5$. For $l\ge k$ both approximation errors are
at most $2^{-r(k)}$, so the right-hand side is at most
$2^{-(k+3)}+2^{-(k+4)}<2^{-k}$. This is stronger than the required
regularity. A represented real threshold is treated by budgeting its rational
approximation error inside the boundary band.
\end{proof}

A locally bounded density supplies the band estimate through
$\mathbb P(t-\delta\le X\le t+\delta)\le2M\delta$ when its density is
bounded by $M$ there. Its finite sampling or transport errors are included
in the certificate. The same lemma can apply to singular continuous laws;
having a density is sufficient, not necessary. Probabilities of hard events
in later statements always require their event names or a stated band
certificate. Finite rational threshold inequalities themselves need no such
extra assumption.

\subsection{Countable operations retain their tail data}

Let $(A_j)$ be a fixed jointly PR family of represented events and let
$H_N=\bigcup_{j<N}A_j$. Suppose a PR function $C$ supplies
\begin{equation}\label{eq:union-cauchy}
 M\ge N\ge C(k)\quad\Longrightarrow\quad
 d_{\rm ev}(H_M,H_N)\le2^{-k}.
\end{equation}
This is a finite-union contract; it makes no reference to an already
constructed infinite union. The regularized finite unions define a
represented event $\bigcup_jA_j$. To obtain an actual cylinder table at
output precision $k$, take $N=C(k+3)$ after monotone closure, approximate
each of its $N$ events with error at most $2^{-(k+3)}/\max(1,N)$, and take
their finite union. The union estimate and \eqref{eq:union-cauchy} give the
required error budget. The schedule is PR; the coarser bound of $N$ additional
guard bits is sufficient if one wishes to avoid logarithms.

A useful sufficient hypothesis is
\[
 M\ge N\ge C(k)\quad\Longrightarrow\quad
 \sum_{j=N}^{M-1}\mathbb P(A_j)\le2^{-k},
\]
by the finite union bound. If the $A_j$ are pairwise disjoint, their finite
union probabilities are the corresponding partial sums; passing through the
same rate proves countable additivity for this presentation. Intersections
follow by complements. This is a substantive countable theorem, but it does
not assert that every countable sequence automatically comes with a tail
modulus.

\subsection{A quantitative convergence theorem without hidden thresholds}

Define $T_K(x)=\max(-K,\min(K,x))$ and
$R_K(x)=(|x|-K)^+$ for rational $K\ge0$. These are continuous, 1-Lipschitz
functions and $|x-T_K(x)|=R_K(x)$.

\begin{lemma}[The clipping estimate]\label{lem:prob-clipping}
For finite rational tables,
\[
 d_1(f,g)\le\max(1,2K)d_0(f,g)+I(R_K(f))+I(R_K(g)).
\]
\end{lemma}
\begin{proof}
Insert $T_K(f)$ and $T_K(g)$ between $f$ and $g$ and use the triangle
inequality. Since clipping is nonexpansive and bounded by $K$,
\[
 |T_K(x)-T_K(y)|\le\min(|x-y|,2K)
                    \le\max(1,2K)\min(1,|x-y|).
\]
Integrate the resulting pointwise inequality. There are no hard completed
threshold events in this proof.
\end{proof}

\begin{theorem}[Effective Vitali completion]
Suppose $f_n$ is a jointly PR family of finite tables with a PR $d_0$ Cauchy
modulus. Suppose also that a PR threshold $K(k)$ satisfies
$I(R_{K(k)}(f_n))\le2^{-k}$ for every $n$. Then the same family has a PR
$d_1$ Cauchy modulus and hence an integrable represented limit.
\end{theorem}
\begin{proof}
For desired error $2^{-k}$ choose $K=K(k+2)$, and choose $h$ with
$\max(1,2K)\le2^h$. Request $d_0(f_m,f_n)\le2^{-(k+h+2)}$ from the
given modulus. Lemma~\ref{lem:prob-clipping} gives
$d_1(f_m,f_n)\le3\cdot2^{-(k+2)}<2^{-k}$. This computes the claimed
modulus, with no extraction from qualitative convergence.
\end{proof}

The same proof gives convergence of an already represented $L^1$ family to
an $L^0$ limit under uniform excess-tail bounds. An $L^1$ majorant supplies
such bounds constructively. If $G\ge0$ has finite approximation $g_r$ and
$K$ bounds every value of $g_r$, then
\[
 \mathbb E(G-K)^+\le\|G-g_r\|_1.
\]
Indeed $(G-K)^+\le|G-g_r|$ in the completed lattice. Therefore
$|X_n|,|X|\le G$ gives uniform tail control for all $X_n$ and $X$.
Combining this with a supplied PR convergence-in-probability rate proves a
quantitative dominated-convergence theorem. This specializes the existing
repository measure framework; neither domination alone nor pointwise
convergence alone is being used to manufacture the missing rate.

\subsection{Conditional expectation on the dyadic filtration}

Let $\mathcal F_n$ denote observation of the first $n$ bits. For a table $f$
at level $m\ge n$, define $C_n f$ by averaging its $2^{m-n}$ entries above
each length-$n$ prefix. If $m<n$, simply lift $f$. Every denominator is a
positive power of two.

\begin{theorem}[Conditional averaging and a quantified martingale limit]
The finite maps $C_n$ preserve positivity and expectation, are $d_1$
contractions, and satisfy $C_nC_m=C_{\min(n,m)}$. They therefore extend to
represented $L^1$ variables. If $X$ has regular name $x_k$, with $x_k$
depending on $\ell(k)$ bits, then
\[
 n\ge\ell(k)\quad\Longrightarrow\quad
 \|C_nX-X\|_1\le2\varepsilon_k.
\]
In particular the sequence $C_nX$ is a represented martingale converging to
$X$ with an explicit PR rate.
\end{theorem}
\begin{proof}
The inequality that the absolute value of an average is at most the average
of the absolute values, summed over prefixes, proves contraction. Regrouping
the same finite sums proves preservation of expectation and the tower law.
For an $\mathcal F_n$ cylinder variable $h$, the same finite regrouping gives
$I(hC_n f)=I(hf)$, the defining conditional-expectation identity. Positivity
is termwise. The completion rule extends these facts to names.

If $n\ge\ell(k)$, then $C_n x_k=x_k$. Hence
\[
 \|C_nX-X\|_1
 \le\|C_n(X-x_k)\|_1+\|x_k-X\|_1
 \le2\|X-x_k\|_1\le2\varepsilon_k.
\]
For error $2^{-p}$, the monotone schedule
$N(p)=\max_{j\le p+1}\ell(j)$ suffices. The tower identity says exactly that
$(C_nX)$ is a martingale. This is convergence of the canonical martingale of
a supplied terminal variable, not an unrestricted theorem for arbitrary
bounded martingales.
\end{proof}

Conditioning on a general represented event $B$ is also available when a
rational margin $\beta>0$ with $\mathbb P(B)\ge\beta$ is supplied:
$\mathbb P(A\mid B)=\mathbb P(A\cap B)/\mathbb P(B)$. Approximate the
denominator enough to keep it above $\beta/2$; ordinary rational division
then has an explicit error budget. A null event supplies no such division,
and a general regular conditional distribution is extra data, not a consequence
of this quotient formula.

\subsection{Independent copies, a uniform law, and a singular continuous law}

A primitive-recursive pairing function partitions the countably many bit
positions into disjoint infinite rows. Copy the table for a represented
variable onto row $i$ by reading only positions assigned to $(i,j)$. Each
finite demand touches finitely many positions and is therefore a cylinder
table; counting unused positions preserves its law and its distance bounds.
The finite product identity proves independence of every finite family of
different rows. Thus a supplied represented random variable has a jointly
PR sequence of independent identically distributed copies. This is an explicit
construction, not an appeal to an arbitrary extension theorem.

Let $U_n$ assign the midpoint $(j+1/2)2^{-n}$ to the $j$th dyadic cell.
For $m\ge n$, $|U_m-U_n|\le2^{-n-1}$ on every common row. The resulting
$U$ is a bounded $L^1$ and $L^0$ variable. Counting cells gives
\[
 \left|\mathbb P(U_n\le t)-\max(0,\min(1,t))\right|\le2^{-n}
\]
for rational $t$, so its CDF is the uniform CDF. The same count supplies
threshold boundary bands of mass at most $2\delta+2^{1-n}$. In particular,
singletons have probability zero; this does not make the value of an actual
infinite sample decidable.

For a continuous law that has no density, use the same bits in
\[
 C_n=\sum_{j=1}^n\frac{2b_j}{3^j}.
\]
The remainder is bounded by $3^{-n}$, so $C_n$ is uniformly Cauchy with an
explicit PR rate. Each prefix of length $n$ puts the limit in one of $2^n$
closed intervals of length $3^{-n}$, with mass $2^{-n}$ each. A singleton
lies in at most one such interval and therefore has probability at most
$2^{-n}$ for every $n$, hence zero. More quantitatively, a sufficiently short
band of width below the level-$n$ gaps meets at most two intervals, giving
a boundary-band modulus and hence a represented continuous CDF.

The same $2^n$ intervals cover the support with total ordinary length
$(2/3)^n$. Choosing $n=3k$ makes this at most $2^{-k}$, since
$(2/3)^3<1/2$. Thus the support is Lebesgue-null with an explicit covering
rate, while the constructed law gives it full mass in the sense that every
finite cover stage has probability one. In standard semantics this is the
Cantor singular continuous distribution. It verifies that the foundation
handles continuous laws beyond absolutely continuous density models.

\subsection{What the arithmetic claim does and does not say}

For each fixed presentation, every displayed finite estimate is an arithmetic
formula over natural codes. The proof recipe consists of exact rational
identities, proved Boolean tests, and quantifier-free induction on finite
partial sums, refinements, and loops. The regularity hypotheses are used at
the finitely many indices demanded by the output precision. A theorem about
all precisions is obtained by universal introduction after those finite
proofs; it is not used as an induction matrix. Completed objects are a
convenient readable notation for this family of contracts.

There is no total universal PR function evaluating every arbitrary PR
program code: if $V(e,n)$ were such a function, then $V(n,n)+1$ would be PR
and its own code would give a contradiction. Accordingly, names here are
fixed PR definitions, possibly with numerical parameters. A variable program
code requires a bounded interpreter and a supplied clock with a proof of
sufficiency. The mathematical claims above are proof-construction schemata
with explicit finite targets; a Lean proof of those targets in ordinary
arithmetic is not yet a term of the embedded object-theory type
$\mathrm{Prf}\;\mathrm{PRAAx}$.

Two elementary obstructions explain why the representation data cannot be
discarded. For a machine index $e$, the stage-$n$ assertion ``$e$ has halted
by stage $n$'' is a PR Boolean sequence, monotone and bounded by one. A
uniform procedure producing its regular real limit to error $1/4$ would
decide halting. Thus bounded monotonicity does not uniformly yield a PR
convergence modulus. Similarly, enumerate no open interval until that machine
halts, and then enumerate $(0,1)$. The resulting effective open set has
uniform probability either zero or one. A uniform regular probability for
all such descriptions would again decide halting. Therefore this report does
not identify all effective Borel descriptions with regular event names.

Exact real equality, arbitrary sigma-algebra completion, automatic rates from
qualitative almost-sure convergence, and general disintegration are not hidden
axioms of the construction. Positive separation, tail bounds, boundary-band
bounds, and convergence rates are retained precisely where the finite proof
uses them. This is what keeps the continuous mathematics within the stated
intuitionistic primitive-recursive proof discipline.
